% Part: lambda-calculus
% Chapter: lambda-definability
% Section: arithmetical-functions

\documentclass[../../../include/open-logic-section]{subfiles}

\begin{document}

\olfileid{lam}{rep}{arf}
\olsection{توابع حسابیِ \usetoken{S}{lambda definable}}

\begin{prop}
  \ollabel{prop:succ-ld}
  تابع جانشین~$\Succ$ !!{lambda definable} است.
\end{prop}

\begin{proof}
ترمی که تابع جانشین را~!!{lambda define}s عبارت است از
\begin{align*}
  \fn{Succ} & \ident \lambd[a][\lambd[fx][f ({a} f x)]].
  \intertext{با قراردادهای ما، این عبارت کوتاه‌شدهٔ عبارت زیر است:}
  \fn{Succ} & \ident \lambd[a][\lambd[f][\lambd[x][(f (({a} f) x))]]].
  \intertext{$\fn{Succ}$ تابعی است که عدد~$a$ را به‌عنوان آرگومان
    می‌پذیرد و به تابع دیگری، یعنی~$\lambd[fx][f ({a} f x)]$، ارزیابی
    می‌شود. آن تابع خود یک عددنمای چرچ نیست. بااین‌حال، اگر آرگومان~$a$
    یک عددنمای چرچ باشد، تابع به یک عددنمای چرچ کاهش می‌یابد. بنگرید:}
   (\lambd[a][\lambd[fx][f ({a} f x)]])\,\num{n} & \redone
   \lambd[fx][f ({\num{n}} f x)].
   \intertext{ترم درونی~$\num{n}fx$ یک سرخ‌عبارت است، زیرا
     $\num{n}$ همان~$\lambd[fx][f^nx]$ است. پس
     $\num{n}fx \red f^nx$ و برای کل ترم داریم:}
   \fn{Succ}\, \num n & \red \lambd[fx][f(f^n(x))],
\end{align*}
یعنی~$\num{n+1}$.
\end{proof}

\begin{ex}
  ببینیم با اعمال~$\fn{Succ}$ بر~$\num{0}$، یعنی~$\lambd[fx][x]$، چه
  رخ می‌دهد. ترم‌ها را به‌طور کامل می‌نویسیم:
  \begin{align*}
    \fn{Succ}\, \num{0} & \ident (\lambd[a][\lambd[f][\lambd[x][(f (({a} f) x))]]])(\lambd[f][\lambd[x][x]])\\
    & \redone \lambd[f][\lambd[x][(f (({(\lambd[f][\lambd[x][x]])} f) x))]] \\
    & \redone \lambd[f][\lambd[x][(f ({(\lambd[x][x])} x))]] \\
    & \redone \lambd[f][\lambd[x][(f x)]] \ident \num{1}\\
  \end{align*}
\end{ex}

\begin{prob}
  ترم
  \begin{align*}
    \fn{Succ}' & \ident \lambd[{n}][\lambd[fx][{n} f (f x)]]
  \end{align*}
  تابع جانشین را~!!{lambda define}s. توضیح دهید چرا.
\end{prob}

\begin{prop}
  \ollabel{prop:add-ld}
  تابع جمع~$\Add$ !!{lambda definable} است.
\end{prop}

\begin{proof}
ترم‌های زیر جمع را~!!{lambda define}:
\begin{align*}
  \fn{Add} & \ident \lambd[{a}{b}][\lambd[fx][{a} f ({b} f x)]]
\intertext{یا، به‌صورت بدیل،}
  \fn{Add}' & \ident \lambd[{a}{b}][{a}\, \fn{Succ} \, {b}].
\intertext{جمع نخست چنین کار می‌کند: $\fn{Add}$ ابتدا دو عدد~${a}$ و
  ${b}$ را می‌پذیرد. حاصل تابعی است که~$f$ و~$x$ را می‌پذیرد و
  $af(bfx)$ را بازمی‌گرداند. اگر~$a$ و~$b$ عددنماهای چرچ~$\num{n}$ و
  $\num{m}$ باشند، این عبارت به~$f^{n+m}(x)$ کاهش می‌یابد، که با
  $f^{n}(f^{m}(x))$ یکسان است. یا، با گام‌های آشکار:}
  (\lambd[{a}{b}][\lambd[fx][{a} f ({b} f x)]])\num n\,\num m & \red
  \lambd[fx][\num{n}\, f (\num {m}\, f x)] \\
  & \red \lambd[fx][\num{n}\, f (f^m x)] \\
  & \red \lambd[fx][f^n (f^m x)] \ident \num {n+m}.
\intertext{بازنمایی دوم جمع، یعنی~$\fn{Add'}$، به شیوه‌ای دیگر کار
  می‌کند. اگر آن را بر دو عددنمای چرچ~$\num{n}$ و~$\num{m}$ اعمال
  کنیم،}
\fn{Add}' \num n \,\num m
& \red \num{n}\, \fn{Succ}\, \num{m}.
\intertext{اما~$\num{n} f x$ همواره به~$f^n(x)$ کاهش می‌یابد. پس}
  \num{n}\,  \fn{Succ}\, \num{m} & \red \fn{Succ}^n(\num{m}).
\end{align*}
و چون~$\fn{Succ}$ تابع جانشین را~!!{lambda define}s و اعمال~$n$بارهٔ
تابع جانشین بر~$m$ برابر~$n+m$ است، این عبارت نیز به~$\num{n+m}$ کاهش
می‌یابد.
\end{proof}

\begin{prop}
  \ollabel{prop:mult-ld}
  ضرب به‌وسیلهٔ ترم زیر~!!{lambda definable} است:
  \[
  \fn{Mult} \ident \lambd[ab][\lambd[fx][a (b f) x]]
  \]
\end{prop}

\begin{proof}
  برای دیدن چگونگی کار این ترم، فرض کنید~$\fn{Mult}$ را بر عددنماهای
  چرچ~$\num{n}$ و~$\num{m}$ اعمال کنیم. ترم
  $\fn{Mult} \, \num{n} \, \num{m}$ به
  $\lambd[fx][\num{n}(\num{m}\, f)x]$ کاهش می‌یابد. ترم~$\num{m} f$
  تابعی را تعریف می‌کند که~$f$ را~$m$ بار بر آرگومانش اعمال می‌کند.
  در نتیجه، $\num{n} (\num{m} f) x$ خود تابع «$f$ را~$m$ بار اعمال
  کن» را~$n$ بار بر~$x$ اعمال می‌کند. به بیان دیگر، $f$ را~$n\cdot m$
  بار بر~$x$ اعمال می‌کنیم. اما ترم نرمال حاصل همان عددنمای چرچ
  $\num{nm}$ است.
\end{proof}

\begin{editorial}
در واقع می‌توانیم با کاهش~$\eta$ این ترم را باز هم ساده‌تر کنیم:
\[
  \fn{Mult} \ident \lambd[ab][\lambd[f][a (b f)]].
\]

اما در آن صورت نخست باید کاهش~$\eta$ را توضیح دهیم.
\end{editorial}

\begin{prob}
ترم زیر ضرب را~!!{lambda define}:
\[
  \fn{Mult}' \ident \lambd[ab][a (\fn{Add}\, b) \num{0}].
\]
توضیح دهید چرا این ترم کار می‌کند.
\end{prob}

تعریف توان‌رسانی به‌صورت یک ترم~$\lambd$ به‌طرز شگفت‌آوری ساده است:
\[
  \fn{Exp} \ident \lambd[be][e b].
\]
آرگومان نخست~$b$ پایه و آرگومان دوم~$e$ توان است. بنا بر کدگذاری اعداد،
به‌طور شهودی~$e f$ همان~$f^e$ است. اگر درک این تعریف دشوار باشد،
توان‌رسانی را می‌توانیم با ضرب تکرارشونده نیز تعریف کنیم:
\[
  \fn{Exp}' \ident \lambd[be][e (\fn{Mult}\, b) \num{1}].
\]

تعریف تابع پیشین و تفریق روی عددنماهای چرچ آن‌قدر که شاید گمان کنیم
ساده نیست؛ این کار به کدگذاری زوج‌ها نیاز دارد.

\end{document}
